\question[4]
Eine Schule setzt ein Informatiksystem ein, in dem Noten, Fehlzeiten und Kontaktdaten der
Schülerinnen und Schüler gespeichert werden. \\
a. Erläutere je zwei technische und zwei organisatorische Maßnahmen, mit denen die Datensicherheit
dieses Systems gewährleistet werden kann. \\
b. Erläutere den Unterschied zwischen \textit{Datensicherheit} und \textit{Datenschutz} an diesem Beispiel.

\begin{solutionbox}{6cm}
a. Technisch z.\,B.: verschlüsselte Speicherung und verschlüsselte Übertragung (HTTPS),
Zugriffskontrolle mit persönlichen Benutzerkonten und starken Passwörtern / Mehrfaktor\-authentifizierung,
regelmäßige Sicherungskopien (Backups), aktuelle Software / Sicherheitsupdates. \\
Organisatorisch z.\,B.: Rollen- und Rechtekonzept (Lehrkraft sieht nur die eigenen Klassen),
Schulung der Mitarbeitenden, klare Regeln zur Weitergabe und Löschung von Daten, protokollierte
Zugriffe. \\
b. \textit{Datensicherheit}: Schutz der Daten vor Verlust, Manipulation und unbefugtem Zugriff
(unabhängig davon, um welche Daten es sich handelt). \textit{Datenschutz}: Schutz des Rechts der
betroffenen Personen auf informationelle Selbstbestimmung -- es dürfen nur die wirklich nötigen
personenbezogenen Daten für einen festgelegten Zweck erhoben und verarbeitet werden.
\end{solutionbox}
